\documentclass[10pt,a4paper]{amsart}
\usepackage[margin=19mm]{geometry}
\usepackage{amsmath,amssymb,mathtools}
\usepackage[scr=rsfs,cal=euler]{mathalfa}
\usepackage{xfrac}
\usepackage[unicode,hidelinks,bookmarks=false]{hyperref}
\newcommand{\ie}{i.e.\ }
\newcommand{\cf}{cf.\ }
\newcommand{\resp}{respectively\ }
\newcommand{\Subsection}{\subsection}
\providecommand{\nobreakdash}{\nobreak\hbox{-}}
\setlength{\emergencystretch}{2em}
\multlinegap0pt
\usepackage{bm}
\usepackage{bbm}
\usepackage{dsfont}
\usepackage{xspace}
\usepackage{cancel}
\usepackage{mathtools}
\usepackage{amsthm}
\makeatletter
\@ifundefined{lemma}{\newtheorem{lemma}{Lemma}}{}
\@ifundefined{proposition}{\newtheorem{proposition}{Proposition}}{}
\makeatother
\def\cN{\mathcal{N}}
\def\to{\rightarrow}
\title[Discovering Causal Relationships using Proxy Variables under]{Discovering Causal Relationships using Proxy Variables under Unmeasured Confounding}
\author{Yong Wu, Yanwei Fu, Shouyan Wang, Yizhou Wang, Xinwei Sun}
\date{Source: arXiv:2510.17167v1 (CC BY 4.0)}
\begin{document}
\maketitle

\noindent\emph{Source and licence.} Discovering Causal Relationships using Proxy Variables under Unmeasured Confounding. Version arXiv:2510.17167v1 (CC BY 4.0), distributed under the Creative Commons Attribution 4.0 International licence, as recorded in the arXiv licence metadata for this exact version and shown on the abstract page https://arxiv.org/abs/2510.17167v1. Full rights notice: licenses/arxiv-2510.17167-CC-BY-4.0.txt.\par\smallskip

\noindent\textit{Editorial scope.} Counted unit: the Proposition whose source label is \texttt{None} in the article (source0.tex lines 3381--3383, source label \texttt{None}), delivered together with the article's own complete original proof of it (source0.tex lines 3384--3491, one \texttt{proof} environment of 108 lines). No mathematics is written, repaired or re-derived: statement and proof are the article's own text line for line. Required context retained uncounted: eq.bridge\_Y (lines 653--655); lemma.miao-example (lines 3977--3984). Every definition, display and label of those lines is retained unchanged. Omitted from this package and not redistributed: 1--652, 656--3380, 3492--3976, 3985--4388. Nothing inside a retained excerpt is omitted or altered: every retained body is delivered exactly as the article's lines are written, so no omission receipt of any kind accompanies this package. Citations: the retained excerpts carry no citation command at all, so no bibliography entry of the article is transcribed and no reference is redistributed. Reference adaptation: none was needed and none was made; every reference inside the delivered unit resolves inside it, so no unresolved reference or citation is produced. Printed slips kept literal: no printed word, symbol, hypothesis or number of the counted statement or of its proof was altered. Layout and class: the article's own class is \texttt{article} with packages including iftex, fontenc, inputenc, textcomp, lmodern, upquote, microtype, parskip, xcolor, longtable, booktabs, array, calc, etoolbox; the wrapper is the lane's pinned \texttt{amsart} editorial wrapper at 10pt on A4, which restates the theorem-like environments the retained text needs and adds the article's own macro definitions that the retained text needs (\texttt{\textbackslash cN}, \texttt{\textbackslash to}), with extra wrapper packages (bm, bbm, dsfont, xspace, cancel, mathtools). The delivered numbering is the unit's own. Assets: no retained span contains a figure environment, a graphics-inclusion command, a PDF-page inclusion, a table or a TikZ picture; no image file is copied into this package, the package declares no asset dependency and no image file is redistributed with it. Licence: version arXiv:2510.17167v1 of this article is distributed under the Creative Commons Attribution 4.0 International licence, as recorded in the arXiv licence metadata for this exact version and shown on the abstract page https://arxiv.org/abs/2510.17167v1. A full rights notice is delivered at licenses/arxiv-2510.17167-CC-BY-4.0.txt. Characters: every retained excerpt is pure ASCII, so the delivered engine, XeLaTeX with its default Latin Modern text font, reports no missing character.
\begin{equation}\label{eq.bridge_Y}
    p(y|x) = \int h(w,y) p(w|x) dw.
\end{equation}

\begin{lemma}
\label{lemma.miao-example}
 If $W|X\sim \cN(\beta_0+\beta_1X,\sigma_{2}^{2})$ and $Y|X\sim \cN(\gamma_0+\gamma_1X,\sigma_{3}^{2})$, then one can verify integral equation $p(y|x)=\int h(w,y)p(w|x)dw$ has a unique solution $h(w,y)$:
\begin{equation}\label{eq.true_solution}
    h(w,y)=\frac{1}{\sigma_{wx}}\phi \left( \frac{y-\gamma_{wx}-\gamma_1/\beta_1 w}{\sigma_{wx}} \right),
\end{equation}
where $\phi$ is the probability density function (pdf) of the standard normal distribution, $\gamma_{wx}=\gamma_{0}-\gamma_{1}\beta_{0}/\beta_{1}$ and $\sigma_{wx}^2=\sigma_3^2-\gamma_1^2\sigma_2^2/\beta_1^2$.
\end{lemma}

\begin{proposition}
Suppose that $X, Y, U, W$ satisfy the linear Gaussian model, \emph{i.e.} $U=\varepsilon_U,X =\alpha_U U + \alpha_0+\varepsilon_X,W =\beta_U U +\beta_0 + \varepsilon_W, Y=\gamma_U U +\gamma_X X+\gamma_W W +\gamma_0 +\varepsilon_Y$, where $\varepsilon_X, \varepsilon_W, \varepsilon_Y, \varepsilon_U$ are standard normal. When $\gamma_W = 0$, as long as $|\gamma_X|>\frac{|B|+\sqrt{\Delta}}{2A}$, where $A=1+\frac{1}{\alpha_{U}^{2}}+\frac{2}{\beta_{U}^{2}}+\frac{1}{\alpha_{U}^{2}\beta_{U}^{2}}+\frac{\alpha_{U}^{2}}{\beta_{U}^{2}}$, $B=\frac{2\gamma_U}{\alpha_U}+\frac{2\gamma_U}{\alpha_U\beta_{U}^{2}}+\frac{2\alpha_U\gamma_U}{\beta_{U}^{2}}$ and $\Delta=\frac{4\left( 1+\alpha_{U}^{2}+\beta_{U}^{2} \right) \left( 1+\alpha_{U}^{2}+\gamma_{U}^{2} \right)}{\alpha_{U}^{2}\beta_{U}^{2}}$, the integration equation \eqref{eq.bridge_Y} has no solution. When $\gamma_W \ne 0$, as long as $|\gamma_W|>\frac{|C|+|B||\gamma_X|+A\gamma_X^2}{2|D|}$, where $C=1-\gamma_U^2/\beta_U^2$ and $D=\frac{\gamma_X}{\alpha_U\beta_U}+\frac{\alpha_U}{\beta_U}\gamma_X+\frac{\beta_U}{\alpha_U}\gamma_X+\frac{\gamma_U}{\beta_U}$, \eqref{eq.bridge_Y} has a solution. 
\end{proposition}

\begin{proof}
Based on the data generation structure, we can obtain the joint distribution
$$
\left( \begin{array}{c}
	U\\
	X\\
	W\\
	Y\\
\end{array} \right) \sim \cN\left\{ \left( \begin{array}{c}
	0\\
	\alpha_0\\
	\beta_0\\
	\gamma_0\\
\end{array} \right) ,\left( \begin{matrix}
	1&		\alpha_U&		\beta_{U}&		\mathrm{Cov}(U,Y)\\
	\alpha_U&		1+\alpha_U^{2}&		\alpha_U\beta_{U}&		\mathrm{Cov}(X,Y)\\
	\beta_{U}&		\alpha_U\beta_{U}&		1+\beta_{U}^{2}&		\mathrm{Cov}(W,Y)\\
	\mathrm{Cov}(U,Y)&		\mathrm{Cov}(X,Y)&		\mathrm{Cov}(W,Y)&		\mathrm{Var}(Y)
\end{matrix} \right) \right\},
$$
where covariance $\mathrm{Cov}(U,Y),\mathrm{Cov}(X,Y),\mathrm{Cov}(W,Y)$ and $\mathrm{Var}(Y)$ are respectively
$$
\left\{ \begin{aligned}
	\mathrm{Cov}(U,Y)&=\gamma_{U}+\gamma_X\alpha_U+\gamma_W\beta_{U}\\
	\mathrm{Cov}(X,Y)&=\alpha_U\left( \gamma_{U}+\gamma_W\beta_{U}+\gamma_X\alpha_U \right) +\gamma_X\\
	\mathrm{Cov}(W,Y)&=\beta_{U}\left( \gamma_{U}+\alpha_U\gamma_X+\gamma_W\beta_{U} \right) +\gamma_W\\
	\mathrm{Var}(Y)&=(\gamma_{U}+\gamma_X\alpha_U+\gamma_W\beta_{U})^2+\gamma_X^{2}+\gamma_W^{2}+1.
\end{aligned} \right.
$$
We can therefore derive the explicit form of the conditional distributions $p(w|x)$ and $p(y|x)$:
$$
\begin{aligned}
	W|X=x&\sim \cN\left\{ \mu_W+\frac{\mathrm{Cov}(W,X)}{\mathrm{Var}(X)}(x-\mu_{X}),\mathrm{Var}(W)\left( 1-\frac{\mathrm{Cov}^2(W,X)}{\mathrm{Var}(X)\cdot \mathrm{Var}(W)} \right) \right\}\\
	&\sim \cN\left\{ \mu_{X}^{W|X}x+\mu_{0}^{W|X},\sigma_{W|X}^{2} \right\} 
\end{aligned}
$$
$$
\begin{aligned}
    Y|X=x&\sim \cN\left\{ \mu_Y+\frac{\mathrm{Cov}(Y,X)}{\mathrm{Var}(X)}(x-\mu_{X}),\mathrm{Var}(Y)\left( 1-\frac{\mathrm{Cov}^2(Y,X)}{\mathrm{Var}(X)\cdot \mathrm{Var}(Y)} \right) \right\}\nonumber \\
	&\sim \cN\left\{ \mu_{X}^{Y|X}x+\mu_{0}^{Y|X},\sigma_{Y|X}^{2} \right\},
\end{aligned}
$$
where $\mu_{X}^{W|X},\mu_{0}^{W|X},\sigma_{W|X}^{2}, \mu_{X}^{Y|X}, \mu_{0}^{Y|X}$ and $\sigma_{Y|X}^{2}$ are defined as follows
$$
\begin{cases}
	\mu_{X}^{W|X}=\frac{\alpha_U\beta_{U}}{1+\alpha_U^{2}}\\
	\mu_{0}^{W|X}=\beta_0-\frac{\alpha_0\alpha_U\beta_{U}}{1+\alpha_U^{2}}\\
	\sigma_{W|X}^{2}=1+\beta_{U}^{2}-\frac{(\alpha_U\beta_{U})^2}{1+\alpha_U^{2}}\\
	\mu_{X}^{Y|X}=\frac{\alpha_U\left( \gamma_{U}+\gamma_W\beta_{U}+\gamma_X\alpha_U \right) +\gamma_X}{1+\alpha_U^{2}}\\
	\mu_{0}^{Y|X}=\gamma_0-\frac{\alpha_0\alpha_U\left( \gamma_{U}+\gamma_W\beta_{U}+\gamma_X\alpha_U \right) +\alpha_0\gamma_X}{1+\alpha_U^{2}}\\
	\sigma_{Y|X}^{2}=(\gamma_{U}+\gamma_X\alpha_U+\gamma_W\beta_{U})^2+\gamma_X^{2}+\gamma_W^{2}+1-\frac{(\alpha_U\left( \gamma_{U}+\gamma_W\beta_{U}+\gamma_X\alpha_U \right) +\gamma_X)^2}{1+\alpha_U^{2}}.
\end{cases}
$$
By applying Lemma~\ref{lemma.miao-example}, the solution of \eqref{eq.bridge_Y} is given by:
$$
h(w,y)=\frac{1}{\sqrt{\sigma_{Y|X}^{2}-\left(\mu_{X}^{Y|X}\right)^2\sigma_{W|X}^{2}/\left(\mu_{X}^{W|X}\right)^2}}\phi \left( \frac{y-\left( \mu_{0}^{Y|X}-\mu_{X}^{Y|X}\mu_{0}^{W|X}/\mu_{X}^{W|X} \right) -\mu_{X}^{Y|X}/\mu_{X}^{W|X}w}{\sqrt{\sigma_{Y|X}^{2}-\left(\mu_{X}^{Y|X}\right)^2\sigma_{W|X}^{2}/\left(\mu_{X}^{W|X}\right)^2}} \right),
$$
where $\phi$ is the standard normal distribution's probability density function (pdf).

For $h(w, y)$ to be meaningful, we need $\sigma_{Y|X}^{2}-\left(\mu_{X}^{Y|X}\right)^2\sigma_{W|X}^{2}/\left(\mu_{X}^{W|X}\right)^2>0$, which implies
\begin{equation}\label{eq. solution_h_exist}
\begin{aligned}
	1-\frac{\gamma_{U}^{2}}{\beta_{U}^{2}}-\left( \frac{2\gamma_{U}}{\alpha_U}+\frac{2\gamma_{U}}{\alpha_U\beta_{U}^{2}}+\frac{2\alpha_U\gamma_{U}}{\beta_{U}^{2}} \right) \gamma_X-\left( 1+\frac{1}{\alpha_U^{2}}+\frac{2}{\beta_{U}^{2}}+\frac{1}{\alpha_U^{2}\beta_{U}^{2}}+\frac{\alpha_U^{2}}{\beta_{U}^{2}} \right) \gamma_X^{2}\\
	-2\left( \frac{1}{\alpha_U\beta_{U}}+\frac{\alpha_U}{\beta_{U}}+\frac{\beta_{U}}{\alpha_U} \right) \gamma_X\gamma_W-2\frac{\gamma_{U}}{\beta_{U}}\gamma_W&>0.
\end{aligned}
\end{equation}

We discuss the following two cases: \textbf{(i)} $X \to Y$ ($\gamma_X \neq 0$) and $W \not\to Y$ ($\gamma_W = 0$); \textbf{(ii)} $X \to Y$ ($\gamma_X \neq 0$) and $W \not\to Y$ ($\gamma_W = 0$).

\textbf{(i). $\gamma_X\ne0,\gamma_W=0$.}

We first rewrite \eqref{eq. solution_h_exist} as:
$$
\underset{C}{\underbrace{1-\frac{\gamma_{U}^{2}}{\beta_{U}^{2}}}}\underset{B}{\underbrace{-\left( \frac{2\gamma_U}{\alpha_U}+\frac{2\gamma_U}{\alpha_U\beta_{U}^{2}}+\frac{2\alpha_U\gamma_U}{\beta_{U}^{2}} \right) }}\gamma_X\underset{A}{\underbrace{-\left( 1+\frac{1}{\alpha_{U}^{2}}+\frac{2}{\beta_{U}^{2}}+\frac{1}{\alpha_{U}^{2}\beta_{U}^{2}}+\frac{\alpha_{U}^{2}}{\beta_{U}^{2}} \right) }}\gamma_{X}^{2}>0.
$$

Noting that this is a quadratic function, we can get its discriminant 
$$
\Delta:=B^2-4AC=\frac{4\left( 1+\alpha_{U}^{2}+\beta_{U}^{2} \right) \left( 1+\alpha_{U}^{2}+\gamma_{U}^{2} \right)}{\alpha_{U}^{2}\beta_{U}^{2}}>0.
$$
Besides, we can find $ 1+\frac{1}{\alpha_{U}^{2}}+\frac{2}{\beta_{U}^{2}}+\frac{1}{\alpha_{U}^{2}\beta_{U}^{2}}+\frac{\alpha_{U}^{2}}{\beta_{U}^{2}}>0$. Therefore, this is a quadratic function whose discriminant is always positive and opens downward. When $\gamma_X$ satisfies $\frac{-B+\sqrt{\Delta}}{2A}<\gamma_X<\frac{-B-\sqrt{\Delta}}{2A}$, \eqref{eq.bridge_Y} will have a solution. When $\gamma_X\geq\frac{-B-\sqrt{\Delta}}{2A}$ or $\gamma_X\leq\frac{-B+\sqrt{\Delta}}{2A}$, \eqref{eq.bridge_Y} will have no solution. 

Without loss of generality, we consider the case where $\alpha_U$ and $\gamma_U$ have the same sign. First, we can find $ B=-(\frac{2\gamma_U}{\alpha_U}+\frac{2\gamma_U}{\alpha_U\beta_{U}^{2}}+\frac{2\alpha_U\gamma_U}{\beta_{U}^{2}})<0$ since $\beta_U^2>0$. Thus, we have $|-B-\sqrt{\Delta}|<|-B+\sqrt{\Delta}|$. Thus, when $|\gamma_X|>\frac{-B+\sqrt{\Delta}}{2A}$,  \eqref{eq.bridge_Y} will have no solution. If $\alpha_U$ and $\gamma_U$ have the different sign, we have $ B=-(\frac{2\gamma_U}{\alpha_U}+\frac{2\gamma_U}{\alpha_U\beta_{U}^{2}}+\frac{2\alpha_U\gamma_U}{\beta_{U}^{2}})>0$ since $\beta_U^2>0$. Thus, we have $|-B-\sqrt{\Delta}|>|-B+\sqrt{\Delta}|$. Thus, when $|\gamma_X|>\frac{-B-\sqrt{\Delta}}{2A}$,  \eqref{eq.bridge_Y} will have no solution.  

Combining the two cases, we can obtain that as long as $|\gamma_X|>\frac{|B|+\sqrt{\Delta}}{2A}$, the integration equation \eqref{eq.bridge_Y} has no solution.

\textbf{(ii). $\gamma_X\ne0,\gamma_W\ne0$.} 

We consider the case $|\gamma_X|>\frac{-B+\sqrt{\Delta}}{2A}$ under $\alpha_U\gamma_U>0$, since  \eqref{eq.bridge_Y} have no solution. We can rewrite  \eqref{eq. solution_h_exist} as
\begin{align*}
    & 2\left( \frac{\gamma_X}{\alpha_U\beta_U}+\frac{\alpha_U}{\beta_U}\gamma_X+\frac{\beta_U}{\alpha_U}\gamma_X+\frac{\gamma_U}{\beta_U} \right) \gamma_W \\
    & <1-\frac{\gamma_{U}^{2}}{\beta_{U}^{2}}-\left( \frac{2\gamma_U}{\alpha_U}+\frac{2\gamma_U}{\alpha_U\beta_{U}^{2}}+\frac{2\alpha_U\gamma_U}{\beta_{U}^{2}} \right) \gamma_X-\left( 1+\frac{1}{\alpha_{U}^{2}}+\frac{2}{\beta_{U}^{2}}+\frac{1}{\alpha_{U}^{2}\beta_{U}^{2}}+\frac{\alpha_{U}^{2}}{\beta_{U}^{2}} \right) \gamma_{X}^{2}.
\end{align*}
Thus, if $\frac{\gamma_X}{\alpha_U\beta_U}+\frac{\alpha_U}{\beta_U}\gamma_X+\frac{\beta_U}{\alpha_U}\gamma_X+\frac{\gamma_U}{\beta_U}<0$, we can obtain that \eqref{eq.bridge_Y} may still have a solution, as long as
$$
\gamma_W>\frac{1-\frac{\gamma_{U}^{2}}{\beta_{U}^{2}}-\left( \frac{2\gamma_U}{\alpha_U}+\frac{2\gamma_U}{\alpha_U\beta_{U}^{2}}+\frac{2\alpha_U\gamma_U}{\beta_{U}^{2}} \right) \gamma_X-\left( 1+\frac{1}{\alpha_{U}^{2}}+\frac{2}{\beta_{U}^{2}}+\frac{1}{\alpha_{U}^{2}\beta_{U}^{2}}+\frac{\alpha_{U}^{2}}{\beta_{U}^{2}} \right) \gamma_{X}^{2}}{2\left( \frac{\gamma_X}{\alpha_U\beta_U}+\frac{\alpha_U}{\beta_U}\gamma_X+\frac{\beta_U}{\alpha_U}\gamma_X+\frac{\gamma_U}{\beta_U} \right)}.
$$
We find that if $\frac{\gamma_X}{\alpha_U\beta_U}+\frac{\alpha_U}{\beta_U}\gamma_X+\frac{\beta_U}{\alpha_U}\gamma_X+\frac{\gamma_U}{\beta_U}<0$, the right-hand side of the above inequality is positive. That means, as long as $|\gamma_W|$ is sufficiently large, the solution to  \eqref{eq.bridge_Y} will still exist when $|\gamma_X|>\frac{-B+\sqrt{\Delta}}{2A}$.


If $\frac{\gamma_X}{\alpha_U\beta_U}+\frac{\alpha_U}{\beta_U}\gamma_X+\frac{\beta_U}{\alpha_U}\gamma_X+\frac{\gamma_U}{\beta_U}>0$, we can obtain \eqref{eq.bridge_Y} may still have a solution, as long as
$$
\gamma_W<\frac{1-\frac{\gamma_{U}^{2}}{\beta_{U}^{2}}-\left( \frac{2\gamma_U}{\alpha_U}+\frac{2\gamma_U}{\alpha_U\beta_{U}^{2}}+\frac{2\alpha_U\gamma_U}{\beta_{U}^{2}} \right) \gamma_X-\left( 1+\frac{1}{\alpha_{U}^{2}}+\frac{2}{\beta_{U}^{2}}+\frac{1}{\alpha_{U}^{2}\beta_{U}^{2}}+\frac{\alpha_{U}^{2}}{\beta_{U}^{2}} \right) \gamma_{X}^{2}}{2\left( \frac{\gamma_X}{\alpha_U\beta_U}+\frac{\alpha_U}{\beta_U}\gamma_X+\frac{\beta_U}{\alpha_U}\gamma_X+\frac{\gamma_U}{\beta_U} \right)}.
$$
We find that if $\frac{\gamma_X}{\alpha_U\beta_U}+\frac{\alpha_U}{\beta_U}\gamma_X+\frac{\beta_U}{\alpha_U}\gamma_X+\frac{\gamma_U}{\beta_U}>0$, the right-hand side is negative. That also means,  as long as $|\gamma_W|$ is sufficiently large, the solution to \eqref{eq.bridge_Y} will still exist when $|\gamma_X|>\frac{-B+\sqrt{\Delta}}{2A}$.

If $\alpha_U\gamma_U<0$, the proof is similar. Besides, in the above cases, as long as $|\gamma_W| >\frac{|C|+|B||\gamma_X|+A\gamma_X^2}{2|D|}$ with $D:=\frac{\gamma_X}{\alpha_U\beta_U}+\frac{\alpha_U}{\beta_U}\gamma_X+\frac{\beta_U}{\alpha_U}\gamma_X+\frac{\gamma_U}{\beta_U}$,  \eqref{eq.bridge_Y} has a solution. 
\end{proof}

\end{document}
